\chapter[Estado del software]{Estado del software}
\label{cp:Estado-software}

\section*{Resumen}

\section{Sistemas operativos y software de las PC por laboratorio}
En todos los laboratorios las computadoras son de escritorio y se intenta que todas las computadoras usadas por los alumnos tengan el mismo hardware, sistema operativo y, sobre todo, la misma configuración de software. Esto es para que los alumnos puedan trabajar en cualquier computadora del laboratorio sin tener que adaptarse a un nuevo entorno de trabajo. 

Las computadoras usadas por el personal técnico de los laboratorios no necesariamente comparten hardware, pero sí el sistema operativo y la configuración de software básico, además de agregar otro software para el diseño y mantenimiento de los equipos del laboratorio. 

En líneas generales, la instalación de software de las PC es:
\subsection{Laboratorio 1}
\begin{itemize}
    \item \textbf{Windows \todo{completar con el sistema operativo}}
    \item \textbf{Google Chrome / Firefox:} Navegadores.
    \item \textbf{Programas \todo{completar con lo que corresponda}}...
\end{itemize}

\subsection{Laboratorio 2}
\begin{itemize}
    \item \textbf{Windows \todo{completar con el sistema operativo}}
    \item \textbf{Google Chrome / Firefox:} Navegadores.
    \item \textbf{Programas \todo{completar con lo que corresponda}}...
\end{itemize}

\subsection{Laboratorio 3}
\begin{itemize}
    \item \textbf{Windows 11 Pro:} Sistema operativo (se usa sin activar por no tener licencia).
    \item \textbf{Google Chrome / Firefox:} Navegadores.
%   \item \textbf{Visual Studio Code:} IDE bastante usado entre alumnos.
%   \begin{itemize}
%       \item Extensiones para desarrollo en Python, C y Arduino.
%   \end{itemize}
    \item \textbf{Anaconda Python:} Con librerías para desarrollo en Python.
    \begin{itemize}
        \item \textbf{Jupyter Notebook:} Uno de los IDEs más usados.
        \item \textbf{Spyder:} Actualmente el IDE más usado.
        \item \textbf{Pandas}
        \item \textbf{NumPy}
        \item \textbf{Matplotlib}
        \item \textbf{SciPy}
        \item \textbf{PyVisa}
        \item \textbf{PySerial:} PyVisa puede reemplazar a PySerial, pero se mantiene por compatibilidad con proyectos viejos.
        \item \textbf{NI-DAQmx:} Para comunicación con las placas de adquisición de NI y Vernier.
    \end{itemize}
    \item \textbf{Arduino IDE}
    \item \textbf{Programas usados con Arduino para ploteo y guardado de datos.}
    \item \textbf{LTspice:} Software de simulación de circuitos electrónicos.
    \item \textbf{NI Package Manager:} Administrar paquetes de software de National Instruments y NI-VISA.
    \begin{itemize}
        \item \textbf{NI-VISA / NI 488.2:} (En general se instalan juntos). Librería para comunicación con instrumentos de medición.
        \item \textbf{NI MAX:} Software para configurar y administrar instrumentos de medición.
        \item \textbf{NI-DAQmx:} Librería para comunicación con las placas de adquisición de NI y Vernier.
    \end{itemize}
    \item \textbf{MotionDAQ:} Software para comunicación con placas Vernier de adquisición de datos.
    \item \textbf{OpenChoice Desktop:} Alternativa para usar con osciloscopios Tektronix.
    \item (OPCIONAL) \textbf{KiCad:} Diseño de PCB.
\end{itemize}

\subsection{Laboratorio 4}
\begin{itemize}
    \item \textbf{Windows 11 Pro:} Sistema operativo (se usa sin activar por no tener licencia).
    \item \textbf{Google Chrome / Firefox:} Navegadores.
    \item \textbf{Visual Studio Code:} IDE bastante usado entre alumnos.
    \begin{itemize}
        \item Extensiones para desarrollo en Python, C y Arduino.
    \end{itemize}
    \item \textbf{Anaconda Python:} Con librerías para desarrollo en Python.
    \begin{itemize}
        \item \textbf{Jupyter Notebook:} Uno de los IDEs más usados.
        \item \textbf{Spyder:} Actualmente el IDE más usado.
        \item \textbf{Pandas}
        \item \textbf{NumPy}
        \item \textbf{Matplotlib}
        \item \textbf{SciPy}
        \item \textbf{PyVisa}
        \item \textbf{PySerial:} PyVisa puede reemplazar a PySerial, pero se mantiene por compatibilidad con proyectos viejos.
        \item \textbf{NI-DAQmx:} Para comunicación con las placas de adquisición de NI y Vernier.
        \item \textbf{OpenCV (cv2):} Para procesamiento de imágenes.
    \end{itemize}
    \item \textbf{Arduino IDE}
    \item \textbf{LTspice:} Software de simulación de circuitos electrónicos.
    \item \textbf{NI Package Manager:} Administrar paquetes de software de National Instruments y NI-VISA.
    \begin{itemize}
        \item \textbf{NI-VISA / NI 488.2:} Librería para comunicación con instrumentos de medición.
        \item \textbf{NI MAX:} Software para configurar y administrar instrumentos de medición.
        \item \textbf{NI-DAQmx:} Librería para comunicación con las placas de adquisición de NI y Vernier.
    \end{itemize}
    \item \textbf{Thorlabs APT:} Software para comunicación con los equipos de medición de Thorlabs.
    \item TODO: "Programas específicos para manejo de equipos Thorlabs, espectrómetros, pinzas ópticas, medidor de potencia, etc."
    \item \textbf{MotionDAQ:} Software para comunicación con placas Vernier de adquisición de datos.
    \item (OPCIONAL) \textbf{Autodesk Inventor:} Software de diseño 3D CAD/CAM/CAE.
    \item (OPCIONAL, no tenemos licencia) \textbf{SolidWorks:} Software de diseño CAD/CAM/CAE.
    \item (OPCIONAL) \textbf{CircuitMaker:} Diseño de PCB.
    \item (OPCIONAL) \textbf{KiCad:} Diseño básico de PCB.
    \item (OPCIONAL, no tenemos licencia) \textbf{MATLAB:} Software de computación numérica y simulación.
    \item (OPCIONAL, no tenemos licencia) \textbf{Altium Designer}.
    \item (OPCIONAL, no tenemos licencia) \textbf{Cadence}.
    \item (OPCIONAL, no tenemos licencia) \textbf{LabVIEW:} Software de programación gráfica para adquisición de datos y control de instrumentos.
\end{itemize}
